\section{Experimental Setup}
\label{sec:setup}

\subsection{Dataset}
The benchmark corpus consists of 8 raw astronomical images, each a
$1500\times1500$ raster of unsigned 16-bit big-endian pixels (4.5\,MB per
file, $2.25\times10^6$ pixels). The images were acquired by different
ground-based telescopes (VLT, TJO, GTC, INT, JKT, LCO, WHT) and span a
deliberate variety of image types: bias frames (near-constant DC pedestal),
flat fields (smooth illumination gradients), and science frames (stellar
fields, galaxies). This diversity ensures that no single predictor type
dominates the results artificially.

\subsection{Reference Compressors}
We compare against the generic tools most commonly used for scientific data:
\texttt{gzip} (levels 1, 6, 9), \texttt{bzip2} (levels 1, 9),
\texttt{xz} (levels 1, 6, 9), and \texttt{zstd} (levels 1, 3, 9, 19).
All are invoked on the raw byte stream without any pre-processing or
endian conversion. Results from two repository-specific codecs
(\texttt{prism} and \texttt{helix}) are included for context.

\subsection{Evaluation Protocol}
The primary metric is bits per byte of the original file:
\[
  \text{bpp} = \frac{8 \times \text{compressed bytes}}{\text{original bytes}}.
\]
For a 16-bit image, $\text{bpp} = \text{bits per pixel}/2$.
Compression and decompression times are measured as wall-clock elapsed time
for a single run on a single machine; all compressors run sequentially in
the same benchmark process to avoid contention.
Losslessness is verified by byte-exact comparison of the decompressed output
against the original file; any mismatch disqualifies the result.
Binary size constraints from the project specification are satisfied:
compress $\le 5$\,MB, decompress $\le 1$\,MB.
